\chapter{REVISÃO BIBLIOGRÁFICA}
Sistemas embarcados com comunicação \textit{IoT} capazes de monitorar a energia elétrica residencial tornaram-se cada vez mais necessários devido às constantes quedas na distribuição de energia e prejuízos relacionados a dispositivos eletroeletrônicos que estão conectados ao sistema elétrico de distribuição. Nesse contexto, faz-se profícua a implementação de estratégias que promovam confiabilidade e eficiência operacional desses sistemas. Diante disso, o uso do microcontrolador \textit{Raspberry Pi Pico 2 W}, equipado com o \textit{chip} RP2350, aliado aos sensores SCT013-100 e ZMPT101B para medição de corrente e tensão, respectivamente, com o uso de \textit{machine learning} para previsão do consumo de energia, têm se mostrado uma abordagem promissora para o monitoramento e previsão do consumo residencial. A seguir, serão apresentados os principais trabalhos correlatos que embasam a implementação do sistema proposto.

No trabalho de \citeonline{Ferraz2018}, foi desenvolvido um analisador de energia monofásico de baixo custo empregando a plataforma de código aberto, com o intuito de facilitar o monitoramento da eficiência energética e democratizar o acesso a essa tecnologia. O sistema utilizou o microcontrolador ATmega328p da plataforma \textit{Arduino}, sensores de corrente e tensão alternada, e um módulo de cartão SD para armazenamento e posterior análise dos dados. O protótipo obteve grande precisão, apresentando erros relativos inferiores a 2,4\% quando comparados a instrumentos comerciais como multímetros.

Em consonância com o avanço da conectividade, \citeonline{Rio2020} apresentaram a plataforma "e-consumo", voltada ao monitoramento inteligente e individualizado do consumo de energia elétrica em residências. O estudo propôs um \textit{hardware} baseado no módulo ESP8266EX, que possui \textit{Wi-Fi} nativo, aliado ao sensor de corrente não invasivo SCT-013. O sistema demonstrou alta eficiência ao enviar dados continuamente via protocolo seguro \textit{MQTT} para um servidor local operado por um \textit{Raspberry Pi}, oferecendo ao usuário transparência e controle de seus gastos por meio de uma interface \textit{web}.

Já no estudo de \citeonline{Bezerra2020}, o foco foi o desenvolvimento de um sistema embarcado de baixo custo para empoderar e auxiliar consumidores do Grupo de Baixa Tensão na adesão a novas modalidades tarifárias do setor elétrico, com ênfase na Tarifa Branca. Utilizando o microcontrolador \textit{Arduino Nano} em função de sua relação de custo e processamento, o projeto foi submetido a testes práticos com diferentes cargas de iluminação e ar-condicionado. Aplicando a análise estatística de Distribuição t de Student, a autora comprovou que o medidor desenvolvido obteve comportamento de medição estatisticamente similar ao de analisadores de mercado, como equipamentos da \textit{Fluke} e \textit{Hikari}.

Um avanço importante no processamento digital de sinais embarcado foi apresentado por \citeonline{oliveira2022}, que desenvolveu um protótipo de medidor monofásico baseado no microcontrolador ESP32. O autor utilizou a técnica da Transformada Rápida de Fourier (FFT) para analisar os sinais dos sensores SCT-013 e ZMPT101B, permitindo o cálculo preciso do fator de potência e o envio periódico dos dados estruturados em \textit{JSON} via \textit{Wi-Fi} para um servidor remoto.

O estudo feito por \citeonline{Pimentel2024} explorou uma nova fronteira ao aliar a Internet das Coisas (\textit{IoT}) a técnicas de Inteligência Artificial para a predição do consumo energético de dispositivos de comunicação LoRa. A metodologia avaliou o uso de algoritmos de \textit{Machine Learning}, como Regressão Linear, Árvore de Decisão (\textit{Decision Tree}) e Floresta Aleatória (\textit{Random Forest}), para estimar o tempo de vida útil das baterias de placas ESP32 em operação contínua e em modo de hibernação. O método \textit{Decision Tree} alcançou alta precisão, com um coeficiente de determinação superior a 96\% na estimativa da autonomia da bateria.

Um protótipo de medidor inteligente (\textit{smart meter}) proposto por \citeonline{Camargo2025} foi desenvolvido especificamente para o cenário do mercado livre de energia brasileiro, com o intuito de reduzir a dependência de dados centralizados pelas concessionárias. O projeto utilizou o microcontrolador ESP32-WROOM acoplado a uma arquitetura de computação em nuvem \textit{serverless} na plataforma AWS (\textit{Amazon Web Services}), integrando funções \textit{Lambda} e bancos de dados para transmitir os dados de forma periódica via HTTP. Embora a comunicação e medição funcionassem, o dispositivo apresentou dificuldades para bater a meta de 5\% de precisão devido ao ruído térmico presente no circuito de condicionamento analógico.

Conforme \citeonline{Flora2025}, o aumento no uso de cargas não lineares tem intensificado a degradação da Qualidade da Energia Elétrica (QEE), motivando o desenvolvimento de analisadores monofásicos de precisão alinhados às normas do Módulo 8 do PRODIST da ANEEL. Este trabalho destacou-se pela utilização do microcontrolador STM32, de maior poder computacional, associado a um transformador de potencial e uma Bobina de Rogowski para medição de corrente escolha justificada por sua imunidade à saturação magnética e alta linearidade. Em ensaios com bancos de cargas resistivas e indutivas, o protótipo alcançou erros inferiores a 3\% ao calcular tensão, corrente, potências e fator de potência, validando-se frente a instrumentos de altíssima precisão como o \textit{Tektronix} PA4000.

No artigo de \citeonline{Lopes2025}, apresentou-se uma aplicação com foco em sustentabilidade, alinhando a eficiência energética à computação verde e aos Objetivos de Desenvolvimento Sustentável (Agenda 2030) da ONU. O trabalho utilizou a placa \textit{Arduino UNO} unida a módulos sensores para monitorar, registrar e quantificar o consumo real de energia elétrica derivado de processos de simulações computacionais complexas especificamente em bioinformática e sequenciamento genético em \textit{clusters}.

Com base no que foi fundamentado nas revisões bibliográficas, fica evidente destacar a relevância da rápida evolução tecnológica dos sistemas embarcados aplicados à medição de energia elétrica ao longo da última década. As tecnologias migraram do mero registro de dados em cartões locais para a adoção de ambientes residenciais sem fio (\textit{Wi-Fi}/\textit{IoT}), processamento robusto em borda (\textit{Edge Computing}) para análise de distúrbios, uso de nuvem para descentralização e integração de inteligência artificial preditiva. Assim, este trabalho tem como objetivo o desenvolvimento do sistema supracitado, estruturado sobre o microcontrolador \textit{Raspberry Pi Pico 2 W}, \textit{chip} RP2350 e sensores de alta precisão (SCT013-100 e ZMPT101B), empregando \textit{machine learning} para fornecer ao usuário residencial um monitoramento exato, preditivo e altamente confiável diante das falhas e demandas das atuais redes elétricas.